\documentclass[11pt]{article}

% ----------- Packages -----------
\usepackage[numbers]{natbib}
\usepackage[margin=1in]{geometry}
\usepackage{amsmath,amssymb,amsthm,mathtools}
\usepackage{enumitem}
\usepackage{microtype}
\usepackage{algorithm}
\usepackage{algpseudocode}
\usepackage[hidelinks]{hyperref}

% nicer comments
\algrenewcommand{\algorithmiccomment}[1]{\hfill\(\triangleright\) #1}


% ----------- Theorems -----------
\newtheorem{definition}{Definition}
\newtheorem{assumption}{Assumption}
\newtheorem{lemma}{Lemma}
\newtheorem{proposition}{Proposition}
\newtheorem{theorem}{Theorem}
\newtheorem{remark}{Remark}
\newtheorem{corollary}{Corollary}

% ----------- Macros -----------
\newcommand{\R}{\mathbb{R}}
\newcommand{\E}{\mathbb{E}}
\newcommand{\KL}{\mathrm{D}_{\mathrm{KL}}}
\newcommand{\Df}{\mathrm{D}_f}
\newcommand{\Pbb}{\mathbb{P}}
\newcommand{\1}{\mathbf{1}}
\newcommand{\push}{\#}

\title{A Neutral Operator Theory for the Information Geometry of Dynamical Systems\\[2pt]
\large From Philosophy to Foundations to Operations}
\author{~}
\date{\today}

\begin{document}
\maketitle

\begin{abstract}
We assume the empirically charitable stance that a dynamical representation is \emph{faithful} if it reproduces one-step transport. This requires treating geometry (embeddings) and probability (kernels) as \emph{co-primitive}, with the joint edge law $\Gamma = \nu \otimes K$ as the basic object.  

Weak axioms reduce admissible losses through $\Gamma$ to Csiszár $f$-divergences, and minimizing this \emph{information action} yields a chart-invariant, compositional principle. In Gaussian limits, constructivist \emph{least-information action} aligns with least physical action.  

Information geometry equips smooth $f$ with the Fisher--Rao metric. The dual affine structure establishes a reflective equilibrium between abstract model and empirical edge laws. While in machine learning the ‘natural gradient’ is canonical only in the parameter-epistemic sense, in our framework Fisher geometry arises from edge laws themselves. This allows us to reinterpret ‘naturalness’ not as an optimization convenience, but as the unique empirical metric that reconciles law and observation. 

Vanishing edge divergence under the dual connections is then equivalent to measure conjugacy between idealized and observed dynamics, yielding generalized embedding guarantees.  We build on this dual foundation and extend it to finite data. Diagnostic exponents, locality $\eta=2s$ (smoothness) and learning $\zeta=2s/(2s+d)$ (dimension), tie smoothness and dimension to convergence rates, while operational tests lend confidence to whether an empirical representation is predictive and faithful to a hypothesized dynamic.
\end{abstract}

\section{Philosophy}
\noindent\textbf{Primitive object.} The empirical unit is $(\nu,K,D)$ via $\Gamma=\nu\otimes K$ and a divergence.

\medskip
\noindent\textbf{Least-information action (LIA).} Fidelity is quantified by edge divergence between data and model edge laws, optimized over representation and kernel.

\medskip


% =========================================================
% =========================================================
\section{Foundations}
We fix notation and state minimal assumptions before developing the divergence-based geometry and its consequences.

\subsection{Standing assumptions (minimal)}
\begin{assumption}[Stationarity]
The source process $(X_t)_{t\in\mathbb{Z}}$ is a time-homogeneous stochastic process on a 
measurable state space $(X,\mathcal{B})$ admitting an invariant probability measure $\mu$. 
That is, for all measurable $A \in \mathcal{B}$ and all $t \in \mathbb{Z}$,
\[
\mathbb{P}(X_t \in A) = \mu(A).
\]
\end{assumption}

\begin{assumption}[Regularization of divergences]
We assume that joint edge laws admit meaningful comparison under the chosen divergence.  
Concretely, this may be ensured by mutual absolute continuity with respect to a common dominating measure, by working on a finite partition, or by introducing a small noise floor to avoid singular supports.  
These conditions guarantee that $D_f(\Gamma_1\|\Gamma_2)=0$ implies $\Gamma_1=\Gamma_2$ almost everywhere, avoiding pathological singularities.
\end{assumption}

\begin{assumption}[Ideal representation]
We posit the existence of an \emph{idealized process} $(X_t)$ with joint edge law $\Gamma$ on $X\times X$, which serves as a conceptual reference.  
The observed process arises through a measurable observation map $\Psi:X\to Y$, producing an empirical kernel $K$ and edge law $\Gamma_K$ on $Y$.  
This assumption reflects our empirical stance: we do not claim the ideal representation is directly accessible, but use it as a standard against which faithfulness can be assessed.
\end{assumption}

\subsection{Objects}
Let $K(y_0,\mathrm dy_1)$ be a Markov kernel on an observation space $Y$.

\begin{definition}[Source law]
The \emph{source law} $\nu$ is the marginal distribution of the initial
state $Y_0$. It represents the empirical frequency with which states
appear as origins of transport, and provides the weighting under which
the kernel $K$ acts.
\end{definition}

\begin{definition}[Joint edge law]
Let $(Y_0,Y_1)$ denote consecutive states of the observed process. 
Given a source law $\nu$ for $Y_0$ and a transition kernel 
$K : \mathcal{Y} \times \mathcal{B}(\mathcal{Y}) \to [0,1]$, the 
\emph{joint edge law} $\Gamma_K$ is the probability measure on 
$\mathcal{Y} \times \mathcal{Y}$ defined by
\[
  \Gamma_K(A \times B) \;=\; \int_A K(y_0,B)\,\nu(\mathrm dy_0),
  \qquad A,B \in \mathcal{B}(\mathcal{Y}).
\]
Equivalently, $\Gamma_K$ admits the disintegration
\[
  \Gamma_K(\mathrm dy_0,\mathrm dy_1) \;=\; 
    \nu(\mathrm dy_0)\,K(y_0,\mathrm dy_1),
\]
so that $Y_0 \sim \nu$ and, conditional on $Y_0=y_0$, we have 
$Y_1 \sim K(y_0,\cdot)$. Thus $\Gamma_K$ describes the law of a 
single transition, or ``edge’’, of the process.
\end{definition}

\begin{definition}[Edge divergence]
Let $D_f$ be a Csiszár $f$-divergence between probability measures.
Given two transition kernels $K_1,K_2$ on a common state space with
common source law $\nu$, their \emph{edge divergence} is
\[
  D_f^{\nu}(K_1,K_2)
  := D_f\!\big(\Gamma_{K_1}\,\big\|\,\Gamma_{K_2}\big),
\]
where $\Gamma_{K_i} = \nu(\mathrm dy_0)K_i(y_0,\mathrm dy_1)$. In other
words, the divergence between kernels is evaluated by lifting them to
their induced joint edge laws.
\end{definition}

\begin{definition}[Primitive empirical object]
The \emph{primitive empirical object} is the triple
\((\nu,K,D_f)\): a source law $\nu$, a one-step kernel $K$, and an
$f$-divergence $D_f$ acting on the joint edge law $\Gamma_K$. This
triple forms the unit of analysis for empirical representations of
dynamics.
\end{definition}

\noindent
By working with $(\nu,K,D_f)$, we treat geometry (embeddings via $\nu$) 
and probability (kernels $K$) as co-primitive, with divergence $D_f$ 
providing the variational principle that links them.

\begin{definition}[Faithfulness under an observation map]
Given a measurable observation map $\Psi:X \to Y$, define the transported edge law
\[
  \Psi^{\natural}(\Gamma) := (\Psi \times \Psi)^{\natural}\,\Gamma,
\]
where $(\cdot)^{\natural}$ abstracts the choice of pushforward or pullback.  
An empirical edge law $\Gamma_K$ is \emph{faithful via $\Psi$} if
\[
  D_f\!\left(\Gamma_K,\,\Psi^{\natural}(\Gamma)\right) = 0.
\]
Faithfulness thus requires that, under $\Psi$, the empirical edge law coincides with the transported ideal law. 
\end{definition}

\subsection{Irreducibility (no-reduction) lemmata}
We first note that neither geometry alone nor probability alone suffices to decide empirical equivalence.

\begin{lemma}[Geometry-only insufficient]\label{lem:G}
There exist systems with identical smooth maps $\psi:Y\to Y$ but distinct kernels, e.g.\ 
deterministic $K^{(1)}(y,\cdot)=\delta_{\psi(y)}$ and noisy 
$K^{(2)}(y,\cdot)=\mathcal N(\psi(y),\sigma^2 I)$, 
having different edge laws. Hence geometry alone cannot decide empirical equivalence 
of one-step transport.
\end{lemma}

\begin{lemma}[Probability-only insufficient]\label{lem:P}
There exist measure-isomorphic systems that are not topologically or smoothly conjugate. 
Measure isomorphism preserves entropy but not continuity or embedding structure. 
Thus probability alone cannot settle geometric equivalence.
\end{lemma}

\subsection{Axioms for admissible empirical loss}
Let $\mathcal L(\nu,K_1,K_2)$ be a loss used to compare dynamics. We posit admissibility conditions:
\begin{itemize}[leftmargin=1.7em,itemsep=2pt,topsep=2pt]
\item \textbf{A1 (Coordinate invariance)}: For bijective measurable $h$,
$\mathcal L(\nu,K_1,K_2)=\mathcal L(h_\push\nu,h_\push K_1,h_\push K_2)$.
\item \textbf{A2 (Data processing)}: For any Markov $M$,
$\mathcal L(\nu,K_1M,K_2M)\le \mathcal L(\nu,K_1,K_2)$ and 
$\mathcal L(M_\push\nu,MK_1,MK_2)\le \mathcal L(\nu,K_1,K_2)$.
\item \textbf{A3 (Temporal additivity)}: For independent blocks, $\mathcal L$ adds over products.
\item \textbf{A4 (Locality)}: Small changes in joint edge laws yield small changes in $\mathcal L$.
\end{itemize}

\begin{proposition}[Axioms force divergence on edge laws]\label{prop:representation}
Under A1--A4, $\mathcal L$ factors through the joint edge laws and is (up to equivalence) 
an $f$-divergence:
\[
\mathcal L(\nu,K_1,K_2)=D_f\!\big(\Gamma_{K_1}\,\big\|\,\Gamma_{K_2}\big).
\]
\end{proposition}

\subsection{Local geometry}
If $K_\theta$ is a smooth family and $f$ is smooth with $f''(1)>0$, then
\[
D_f^\nu(K_{\theta+\delta\theta},K_\theta)
= \tfrac12\,\delta\theta^\top I_f(\theta)\,\delta\theta
+ o(\|\delta\theta\|^2),
\]
which defines the Fisher--Rao (information) metric $I_f$ on the parameter manifold~$\Theta$. 
This is the intrinsic, reparameterization--invariant geometry relevant for local inference.

\paragraph{State-space pullback.}
Viewing $y\mapsto K(\cdot|y)$ as a map into a statistical manifold, pulling back Fisher yields
\[
G(y) \;:=\; \E_{K(\cdot|y)}\!\big[\nabla_y \log k(Y_1|y)\,\nabla_y \log k(Y_1|y)^\top\big],
\]
a tensor field that measures how distinguishable outgoing conditionals are under small moves in $y$.

In practice we use the stable mean-only proxy 
\[
J(y)\;\approx\;(\partial_y\mu(y))^\top \Sigma(y)^{-1} (\partial_y\mu(y)),
\]
where $\mu(y)=\E[Y_1|y]$ and $\Sigma(y)=\mathrm{Var}[Y_1|y]$. Basis/lag selection is then posed 
as maximizing Fisher capture $\mathrm{trace}(L^\top \bar J L)$ or $\log\det(L^\top \bar J L)$, 
with $\bar J := \E_\nu[J(Y)]$.

\paragraph{Gaussian local model.}
If $K(y,\cdot)\approx \mathcal N(\mu(y),\Sigma(y))$ with smooth $\mu,\Sigma$, neglecting $\partial_y\Sigma$ gives
\[
J(y)\;\approx\;(\partial_y\mu(y))^\top \Sigma(y)^{-1} (\partial_y\mu(y)),
\]
the familiar Gauss--Newton/Fisher form. In a local linear fit $Y_{t+1}\approx A(y_t)Y_t+b(y_t)$,
we set $\partial_y\mu(y)\!\approx\!A(y)$ so $J(y)\!\approx\!A(y)^\top \Sigma(y)^{-1} A(y)$.

\begin{remark}[Minimality of the triple $(\nu,K,D_f)$]
Remove $K$ $\Rightarrow$ no transport/prediction; remove $D_f$ $\Rightarrow$ no ruler/decision; 
remove $\nu$ $\Rightarrow$ no source weighting. Add path laws and additivity collapses them back 
to edge laws. Hence $(\nu,K,D_f)$ is both sufficient and minimal for empirical comparison.
\end{remark}

\begin{remark}[Why one-step is primitive]
Pathwise divergences (e.g.\ KL) decompose into sums of edge-level terms for Markov/i.i.d.\ blocks; 
thus the edge divergence is the atomic unit, with path ``action'' an additive aggregate.
\end{remark}

\subsection{From local to global}
\begin{lemma}[Zero divergence $\Rightarrow$ kernel equality]\label{lem:zero}
If $D_f(\Gamma_{K_1}\|\Gamma_{K_2})=0$ and $\Gamma_{K_1}\ll \Gamma_{K_2}$ (or conversely), 
then $K_1(y_0,\cdot)=K_2(y_0,\cdot)$ for $\nu$-a.e.\ $y_0$.
\end{lemma}

\begin{theorem}[Sufficiency $\Rightarrow$ Fisher equality]\label{thm:fisher}
If $\Psi$ is sufficient for $\Gamma$, then for any smooth Csiszár $f$-divergence with $f''(1)>0$, 
the local quadratic form is preserved:
\[
I_{\Gamma}(\theta) \;=\; I_{\Psi(\Gamma)}(\theta).
\]
\end{theorem}

\begin{remark}[Local vs global]
Sufficiency preserves Fisher information (a local, second-order object), 
but does not by itself imply equality of edge laws. 
Global conjugacy requires vanishing divergence, a stronger condition we develop below.
\end{remark}

\begin{theorem}[Generalized embedding theorem]\label{thm:bicond}
Let $h=\Psi_k$ be an embedding map with $\nu=h_\push\mu$. With $K_X$ the kernel on $X$ and $K_Y$ on $Y$,
\[
D_f^{\nu}\!\big(K_Y,\ h_\push K_X\big)=0
\quad\Longleftrightarrow\quad
\text{$K_X$ and $K_Y$ are measure-conjugate via $h$}.
\]
\end{theorem}
\begin{proof}[Sketch]
Zero divergence and domination imply $K_Y(y,\cdot)=h_\push K_X(y,\cdot)$ for $\nu$-a.e.\ $y$ (Lemma~\ref{lem:zero}). 
Iterating $h$ gives equality of $n$-step cylinder measures, hence measure conjugacy.
\end{proof}

\begin{corollary}[Faithfulness $\Rightarrow$ measure conjugacy]
If $\Gamma_K$ is faithful to $\Gamma$ via a diffeomorphic $\Psi$, i.e.\ $D_f(\Gamma_K,\Psi(\Gamma))=0$, 
then $\Gamma_K$ and $\Gamma$ are measure-conjugate through $\Psi$.
\end{corollary}


% ====== INSERTED: Dual affine geometry now lives in Foundations ======
\subsection{Dual affine geometry of divergence}

The Fisher metric $g$ equips the statistical manifold of kernels with a Riemannian structure,
but curvature alone does not explain why divergences decompose additively or why certain
parameterizations are privileged. A further refinement is provided by \emph{dual affine
connections}, which pair with the Fisher metric to yield a dually flat structure.

\begin{definition}[Dual connections]
Let $(\mathcal M,g)$ be a statistical manifold with Fisher metric $g$.
Two affine connections $\nabla$ and $\nabla^\ast$ are said to be \emph{dual} with respect
to $g$ if for all vector fields $X,Y,Z$,
\[
X[g(Y,Z)] \;=\; g(\nabla_X Y,Z) \;+\; g(Y,\nabla^\ast_X Z).
\]
\end{definition}

Amari's $\alpha$--connections $\{\nabla^{(\alpha)} : \alpha\in\mathbb R\}$ form
a one-parameter family of mutually dual connections. Of particular importance are
$\alpha=1$ (the exponential or $e$--connection) and $\alpha=-1$ (the mixture or $m$--connection).

\begin{remark}[Exponential and mixture flatness]
Exponential families are affine under $\nabla^{(e)}$, so their natural parameters form an
$e$--flat coordinate system. Mixture families are affine under $\nabla^{(m)}$, so their
mixture weights form an $m$--flat coordinate system. The two are Legendre dual via the
log-partition function.
\end{remark}

\begin{definition}[Amari projections]
Given a distribution $p$ and a model family $\mathcal M$,
\begin{align*}
\text{m-projection:} & \quad q^\ast = \operatorname*{arg\,min}_{q\in\mathcal M} D_{\mathrm{KL}}(p\|q), \\
\text{e-projection:} & \quad q^\ast = \arg\min_{q\in\mathcal M} D_{\mathrm{KL}}(q\|p).
\end{align*}
The m-projection is orthogonal with respect to the $m$--connection,
while the e-projection is orthogonal with respect to the $e$--connection.
\end{definition}

\begin{theorem}[Pythagorean theorem for KL divergence]
If $\mathcal M$ is $e$--flat and $\mathcal N$ is $m$--flat, then for $p\in\mathcal N$ and
$q^\ast$ the m-projection of $p$ onto $\mathcal M$,
\[
D_{\mathrm{KL}}(p\|q) \;=\; D_{\mathrm{KL}}(p\|q^\ast) + D_{\mathrm{KL}}(q^\ast\|q),
\qquad \forall q\in\mathcal M.
\]
Thus KL divergence decomposes orthogonally into increments along the dual affine flats.
\end{theorem}

\begin{corollary}[Edge laws as dual projections]
Let $\Gamma_{\text{data}}$ denote the empirical edge law, and let
$\mathcal M=\{\Gamma_K : K \in \mathcal K\}$ be an $e$--flat family of candidate kernels.
Then the $m$--projection of $\Gamma_{\text{data}}$ onto $\mathcal M$ is the model kernel
$\Gamma_{K^\ast}$ minimizing the divergence,
\[
\Gamma_{K^\ast} = \arg\min_{\Gamma_K \in \mathcal M} D_{\mathrm{KL}}(\Gamma_{\text{data}}\|\Gamma_K).
\]
For any $\Gamma_K \in \mathcal M$,
\[
D_{\mathrm{KL}}(\Gamma_{\text{data}}\|\Gamma_K)
= D_{\mathrm{KL}}(\Gamma_{\text{data}}\|\Gamma_{K^\ast})
+ D_{\mathrm{KL}}(\Gamma_{K^\ast}\|\Gamma_K).
\]
Thus the excess divergence of any model candidate decomposes into the irreducible divergence
between data and its $m$--projection, plus the divergence between this projection and the candidate.
\end{corollary}

\begin{corollary}[General $f$--divergences]
For a smooth $f$--divergence with $f''(1)>0$, the exact Pythagorean relation does not hold,
but the local quadratic approximation is preserved. For $q^\ast$ the $m$--projection of
$p$ onto $\mathcal M$,
\[
D_f(p\|q) \;=\; D_f(p\|q^\ast) + \tfrac{1}{2}\,\delta\theta^\top I_f(\theta^\ast)\,\delta\theta
\;+\; o(\|\delta\theta\|^2),
\]
where $\delta\theta$ is the displacement from $q^\ast$ to $q$ in parameter space and
$I_f$ is the Fisher information metric induced by $D_f$. Thus the decomposition holds
up to second order, with increments governed by Fisher geometry. \emph{In particular, the
quadratic term controls the growth rates of local divergence, and hence directly determines
the locality and learning exponents $(\eta,\zeta)$ introduced earlier.}
\end{corollary}

\begin{proposition}[Rose condition as intersection of dual flats]
Empirical edge laws, constructed as mixtures of observed conditionals, live on an $m$--flat
submanifold. Candidate kernels, parameterized as local exponential families, live on an
$e$--flat submanifold. The Rose condition
\[
D_f(\Gamma_{\text{data}},\Gamma_{K}) = 0
\]
is equivalent to $\Gamma_{\text{data}}$ lying in the intersection of the $m$--flat and
$e$--flat submanifolds. In this case, the $m$--projection coincides with the model itself,
and the divergence vanishes.
\end{proposition}
% ====== END INSERT ======


\subsection{Sufficient statistics as observation functions}

Chentsov's theorem identifies the Fisher metric as the unique (up to rescaling) 
Riemannian metric monotone under Markov morphisms (statistics). 
Amari's $\alpha$--connections $\{\nabla^{(\alpha)}\}$ preserve this structure, 
and when a statistic is sufficient the induced Fisher metric and dual connections 
are isometric to the originals. Thus sufficiency is not a lossy compression but a 
geometry-preserving reduction.

\medskip

\noindent\textbf{Conditional families as statistical manifolds.}
In our setting the one-step conditionals 
\[
\mathcal F \;=\;\{\,K(\cdot\mid y): y\in Y\,\}
\]
form a statistical manifold, with the state $y$ acting as a parameter. 
If these conditionals are locally exponential families with sufficient statistic 
$T$, then the expectation parameter
\[
m(y) \;:=\; \mathbb E[T(Y_1)\mid Y_0=y]
\]
provides an $m$--affine coordinate system on~$\mathcal F$.

\begin{lemma}[Expectation map as embedding]
Suppose $m:A\to\mathbb R^k$ is $C^1$ on the attractor $A\subseteq Y$, 
with $\mathrm{rank}\,\partial m(y)=\dim A$ for all $y\in A$ (immersion), 
and $m$ is injective on $A$. Then $F(y)=m(y)$ is a diffeomorphic embedding of 
$A$ onto its image. In this case, the sufficient statistic $T$ both preserves the 
$\alpha$--geometry of the conditional family and yields an observation function 
that separates states.
\end{lemma}

\begin{proof}[Sketch]
Immersion + injectivity ensures $m$ provides smooth coordinates 
on the attractor, hence a diffeomorphic embedding.
\end{proof}

\begin{corollary}[Gaussian case]
If $K(\cdot\mid y)\approx \mathcal N(\mu(y),\Sigma)$ with fixed $\Sigma$, then 
$T(y_1)=(y_1,y_1y_1^\top)$ is sufficient and 
$m(y)=(\mu(y),\,\mu(y)\mu(y)^\top+\Sigma)$. 
If $\mu:Y\to\mathbb R^d$ is an injective immersion on $A$, then $F(y)=\mu(y)$ 
itself is an embedding. This coincides with the familiar local Fisher picture, 
where $J(y)\approx (\partial_y\mu(y))^\top\Sigma^{-1}(\partial_y\mu(y))$.
\end{corollary}

\medskip

Thus sufficiency on the model manifold and embedding on the state manifold 
can coincide: the same statistic that preserves Fisher geometry of the 
conditional family can simultaneously serve as a Takens-style observation map. 
The speed at which such embeddings and sufficiency conditions become empirically 
detectable is governed by divergence decay rates, linking them directly to 
ergodic exponents and the finite-data diagnostics $(\eta,\zeta)$ introduced next.


% ====== MOVED HERE: Rates section after Dual affine geometry ======
\subsection{Rates of divergence and dynamical critical exponents}

The asymptotic decay of divergences under e- and m-projections reflects the same scaling laws that ergodic theory encodes in its critical exponents. In the limit $n,t\to\infty$, both projections converge to zero if the empirical law $\Gamma_K$ is faithful to the embedded truth $\Psi(\Gamma)$, but the \emph{rates} of convergence reveal which features of the dynamics are bulk-driven (typical set) and which are tail-driven (rare events). Here the slopes $(\zeta,\eta)$ operationalize these rates: $\eta$ reflects local Fisher curvature, while $\zeta$ interpolates smoothness and effective dimension.

\paragraph{Dictionary.}
\begin{center}
\begin{tabular}{|l|l|l|}
\hline
\textbf{Exponent} & \textbf{Dynamics meaning} & \textbf{Divergence signal} \\
\hline
Entropy $h$ & Orbit complexity growth & Rate of per-step KL (cross-entropy) \\
Lyapunov $\lambda$ & Stretching of trajectories & Tail sensitivity, rare-event amplification \\
Recurrence $\rho$ & Return-time scaling & Limits e-projection convergence (tails) \\
Dimension $d_\mu$ & Mass scaling/fractality & Gap between e- and m-projection rates \\
Mixing $\kappa$ & Decay of correlations & Divergence contraction $\sim t^{-\kappa}$ \\
Spectral $\beta$ & Transfer/Koopman decay & Exponential decay via spectral gap \\
\hline
\end{tabular}
\end{center}

% Defer detailed lemmas/proofs to the appendix
\paragraph{Bridge theorem (informal).}
Classical ergodic exponents (entropy, Lyapunov, recurrence, dimension, mixing, spectrum)
manifest directly as rates at which divergences vanish.
Proof sketches and detailed statements are collected in the appendix.
% ====== END MOVE ======

% =========================================================
\section{Finite data: convergence \& irreducible randomness}

\paragraph{Two asymptotic lenses (empirical diagnostics).}
Let \(\widehat{\mathcal A}\) denote the estimated edge action on held-out data.
\begin{itemize}[leftmargin=1.7em,itemsep=1.5pt,topsep=2pt]
\item \textbf{Learning curve (sample size)}: with model capacity tuned by CV at each size \(n\),
\[
\widehat{\mathcal A}(n) \;\approx\; \mathcal A_\infty \;+\; C_{\mathrm{learn}}\, n_{\mathrm{eff}}^{-\zeta}, \quad n_{\mathrm{eff}}:=n/\tau_{\mathrm{mix}}.
\]
\item \textbf{Locality curve (bandwidth)}: for large \(n\), with a local estimator of scale \(h\),
\[
\widehat{\mathcal A}(h) \;\approx\; \mathcal A_0 \;+\; C_{\mathrm{loc}}\, h^{\eta}.
\]
\end{itemize}
Here \(\tau_{\mathrm{mix}}\) is an integrated autocorrelation time (effective sample-size correction). The asymptotes \(\mathcal A_\infty\) and \(\mathcal A_0\) coincide (up to measurement-noise corrections) in well-specified settings.

\begin{proposition}[Bias--variance law for excess edge action]
For a Markov system with Hölder smooth conditional of order \(s>0\), a local polynomial estimator of order matching \(s\), embedding dimension \(d\), and effective size \(n_{\mathrm{eff}}\), the excess action admits the decomposition
\[
\E\big[\widehat{\mathcal A}(h,n)\big]-\mathcal A_0
\;\approx\;
C_b\, h^{2s} \;+\; C_v\,\frac{1}{n_{\mathrm{eff}}\,h^{d}},
\]
where \(C_b\) (curvature/Fisher-weighted) and \(C_v\) (noise/mixing-dependent) are positive constants.
\end{proposition}

\begin{corollary}[Linking exponents to smoothness and dimension]
If variance is negligible (large \(n\)), then \(\widehat{\mathcal A}(h)-\mathcal A_0\propto h^{\eta}\) with
\[
\boxed{\ \eta=2s\ }.
\]
Balancing bias and variance gives the optimal scale \(h^\star\asymp n_{\mathrm{eff}}^{-1/(2s+d)}\) and the learning rate
\[
\boxed{\ \zeta=\frac{2s}{2s+d}\ }.
\]
Thus \(\eta\) reveals smoothness \(s\approx \eta/2\), and with \(\zeta\) one can back out an \emph{effective} dimension
\[
\boxed{\ d \approx \frac{2s(1-\zeta)}{\zeta}\ }.
\]
\end{corollary}

\begin{remark}[Stochastic floors and irreducible randomness]
If (after correcting for measurement noise) both asymptotes satisfy
\(\ \mathcal A_\infty \approx \mathcal A_0 > 0\), then no further geometric refinement (embedding), capacity growth, or locality contraction will collapse the action: this is a \emph{signature of irreducible process noise}. If either asymptote tends to \(0\), the apparent randomness was due to aliasing, underfitting, or finite data.
\end{remark}

% =========================================================
\section{Pathwise and deterministic limits (how least physical action emerges)}
We recover least physical action as a continuum limit of summed edge actions under forward-KL and Gaussian local models.
\begin{proposition}[Gaussian forward-KL $\Rightarrow$ quadratic path action]
Let time be discretized with step \(\Delta t\). Suppose the model kernel is locally Gaussian:
\[
K_\theta(y,dy') \approx \mathcal N\!\big(y+\mu_\theta(y)\,\Delta t,\ \Sigma(y)\,\Delta t\big).
\]
Then the summed edge action \(\sum_t \KL\big(\Gamma_{\text{data}}^{(t)}\ \|\ \Gamma_{K_\theta}^{(t)}\big)\)
Riemann--sums to a quadratic path functional
\[
\mathcal A_{\mathrm{path}}(\theta)
= \tfrac12\,\E\!\int_0^T \big\|\dot y_t-\mu_\theta(y_t)\big\|^2_{\Sigma(y_t)^{-1}}\,dt\ +\ \text{const}.
\]
\end{proposition}

\begin{remark}[Onsager--Machlup/Girsanov bridge]
For diffusions \(dY_t=b(Y_t)\,dt+\sigma\,dW_t\), changing drift to \(b+u\) yields a path-space KL equal to
\(\tfrac12\E\!\int_0^T \|u_t\|^2_{(\sigma\sigma^\top)^{-1}}\,dt\). Thus most-probable paths under small noise extremize the quadratic stochastic action, aligning LIA with Onsager--Machlup.
\end{remark}

\begin{proposition}[Small-noise/continuum limit $\Rightarrow$ least physical action]
In the limit of vanishing noise and \(\Delta t\to 0\), stationary paths of \(\mathcal A_{\mathrm{path}}\) satisfy Euler--Lagrange/Hamilton--Jacobi equations associated with the induced quadratic Lagrangian. Hence classical least action is recovered as a limit of LIA.
\end{proposition}

% =========================================================
\section{Operations (how it collapses practice)}
This section collapses the theory into a single scalar diagnostic (edge action) with two slopes $(\zeta,\eta)$ and an information-guided selector.

\paragraph{One scalar: edge action.}
Fix a representation and divergence; estimate on a held-out block the \emph{edge action}
\[
\mathcal A\ :=\ \widehat{\Df}^{\nu}\!\big(K_{\text{data}},\ K_{\text{model}}\big),
\]
with a block-bootstrap confidence interval. Decision: CI contains \(0\) (and mean tiny) \(\Rightarrow\) empirical Rose pass; with Takens, invoke Theorem~\ref{thm:bicond}. Else, adjust basis or kernel.

\paragraph{Optional $e$--side diagnostic (capacity sweep).}
To expose the complementary $e$–side contraction, construct a capacity/regularization sweep
for the model family $\{\Gamma_\theta\}$ and track the \emph{$e$–excess}
\(
D_{\mathrm{KL}}(\Gamma_\theta\Vert \Gamma_{\text{data}}) - D_{\mathrm{KL}}(\Gamma_{\theta^\ast}\Vert \Gamma_{\text{data}})
\).
On log–log axes (excess vs.\ effective capacity or sample size), the slope estimates the $e$–side
learning exponent $\zeta$. In practice this uses the same kernels/embeddings as the $m$–side, and
shares hyperparameter infrastructure (CV/bandwidth grids), so it adds a single loop without
altering the main workflow.

\paragraph{Two robust implementations (pick one).}
\begin{itemize}[leftmargin=1.7em,itemsep=2pt,topsep=2pt]
\item \textbf{VQ--JS (discrete, stable):} global codebook on \(y_1\) (k-means, \(M\) cells). For each held-out \(y_0\):
empirical conditional (kNN over edges) \(\to\) histogram \(p\) (smoothed);
model conditional (sample from local Gaussian kernel) \(\to\) histogram \(q\) (same cells, smoothed);
compute \(\text{JS}(p,q)\) and average across \(y_0\).
\item \textbf{Parzen--KL (continuous, MC):} smooth empirical conditional \(\hat p_b(y_1|y_0)\) (Parzen); approximate \(D_{\mathrm{KL}}(\hat p_b\|\hat K)\) by Monte Carlo; average across \(y_0\).
\end{itemize}

\paragraph{Kernel estimator (model).}
Local linear mean via weighted least squares \(Y_{t+1}\approx A(y_0)Y_t+b(y_0)\), residual diffusion
\[
\Sigma(y_0)=\frac{\sum w_t r_t r_t^\top}{\sum w_t}+\lambda I,\quad r_t=Y_{t+1}-A(y_0)Y_t-b(y_0),
\]
and \(\hat K(y_0,\cdot)=\mathcal N(\mu(y_0),\Sigma(y_0))\) (mixtures as needed).
Tune neighborhood scale \(h\), ridge \(\lambda\) (and mixture size) by held-out conditional log-likelihood.

\paragraph{Information-guided basis and lag selection.}
Aggregate the pullback metric by visitation to obtain a global selector
\[
\bar J\ :=\ \E_\nu[J(Y)]\ \approx\ \frac{\sum_t w_t\,A_t^\top\,\Sigma_t^{-1}\,A_t}{\sum_t w_t},
\]
using the same locality weights $w_t$ as in the kernel fit.

\emph{Basis (rotation).} For a linear projection $L\in\R^{d\times r}$ with $L^\top L=I_r$, maximize the
\emph{Fisher capture} $\mathrm{trace}(L^\top \bar J L)$. The optimizer is given by the top-$r$
eigenvectors of $\bar J$ (information-PCA). This choice is invariant to invertible linear reparameterizations.

\emph{Lag (subset) selection.} When $y$ is a delay stack, $\bar J$ is block-structured by lags.
Two principled criteria are:
\begin{itemize}[leftmargin=1.4em,itemsep=1pt,topsep=2pt]
\item \textbf{A-optimal (trace):} choose $S$ to maximize $\mathrm{trace}(\bar J_{SS})$;
\item \textbf{D-optimal (log-det):} choose $S$ to maximize $\log\det(\bar J_{SS})$ (encourages non-redundant lags).
\end{itemize}
A stable greedy D-optimal procedure:
\begin{enumerate}[leftmargin=1.4em,itemsep=1pt,topsep=2pt]
\item Initialize $S\gets\varnothing$.
\item For each candidate $\ell\notin S$, compute the Schur-complement gain
$\Delta_\ell=\log\det\!\big(J_{\ell\ell}-J_{\ell S}J_{SS}^{-1}J_{S\ell}\big)$.
\item Add $\ell^\star=\arg\max_\ell \Delta_\ell$ to $S$ and iterate until $\Delta_{\ell^\star}$ is below a tolerance or a CV-selected budget is reached.
\end{enumerate}

\emph{Refit \& certify.} After selecting $L$ (and/or $S$), refit the kernel in the new representation and
recompute (i) the edge action with a block-bootstrap CI and (ii) the learning/locality exponents $(\zeta,\eta)$.
A better basis typically lowers the edge action (or tightens its CI), increases $\zeta$, and steepens~$\eta$.

\emph{Notes.} Pre-whiten $y$ for stability; use a small ridge on $\Sigma_t$.
If $\Sigma(y)$ varies strongly, include the covariance-variation contribution to $J(y)$; otherwise the mean term dominates in practice.


\paragraph{Outputs.}
Report \(\mathcal A\) with CI (primary), predictive NLL (sanity), and optionally moment gaps / quick spectral snapshot. 

% =========================================================
\section{Checklist (logical closure)}
\begin{itemize}[leftmargin=2em,itemsep=2pt,topsep=2pt]
  \item \textbf{Irreducibility:} Geometry or probability alone is insufficient (Lemmas~\ref{lem:G}, \ref{lem:P}).
  \item \textbf{Axioms $\Rightarrow f$-divergences:} Proposition~\ref{prop:representation}.
  \item \textbf{Zero-divergence lemma:} Equality-from-zero (Lemma~\ref{lem:zero}).
  \item \textbf{Takens+embedding:} With a delay embedding, vanishing edge divergence implies measure conjugacy (Theorem~\ref{thm:bicond}).
  \item \textbf{Dual-affine geometry:} Projections and Pythagorean decompositions explain additive divergence structure.
  \item \textbf{Finite-data diagnostics:} Exponents $(\eta,\zeta)$ expose smoothness and dimension, and operational tests certify faithfulness.
  \item \textbf{Continuum link:} Gaussian/forward-KL limits recover least physical action.
\end{itemize}


\nocite{*}
\bibliographystyle{plainnat}
\bibliography{lia_refs}
\appendix

\section{Divergences and Critical Exponents}
\label{app:exponents}

This section collects the technical lemmas linking divergence decay rates to classical
dynamical critical exponents.

\begin{lemma}[Contraction under transfer operator]
Let $P$ be the Perron–Frobenius operator with invariant measure $\nu$. For any $f$-divergence,
\[
D_f(P\mu \Vert P\nu) \le \eta_f(P)\, D_f(\mu \Vert \nu),
\]
where $\eta_f(P)\in[0,1]$ is a contraction coefficient. If $P$ has a spectral gap with decay $e^{-\gamma t}$, then $D_f(P^t\mu\Vert P^t\nu)$ decays at rate $\gamma$, corresponding to the system's mixing/spectral exponent.
\end{lemma}

\begin{lemma}[Entropy rate]
For stationary processes,
\[
\frac{1}{n} D_{\mathrm{KL}}(\Gamma^{\otimes n}\Vert \Gamma_K^{\otimes n}) \to H(\Gamma \Vert \Gamma_K),
\]
the cross-entropy rate. If $\Gamma=\Gamma_K$, the limit vanishes and the slope constants are controlled by the entropy rate $h$.
\end{lemma}

\begin{lemma}[Recurrence tails]
If return times $\tau_\varepsilon$ satisfy $\mathbb P(\tau_\varepsilon > t) \sim t^{-\rho}$, then under-coverage of rare returns limits the rate of e-projection convergence by $\rho$.
\end{lemma}

\begin{lemma}[Dimension exponent]
Locality and learning curves obey
\[
\hat A(h,n)-A_0 \asymp C_b h^{2s} + \frac{C_v}{n_{\mathrm{eff}} h^{d_\mu}},
\]
so the effective dimension $d_\mu$ appears as the scaling exponent in divergence decay. The gap between e- and m-projections is larger when $d_\mu$ is small (singular measures).
\end{lemma}

\begin{theorem}[Bridge between divergences and exponents]
Let $P$ be the transfer operator of a stationary dynamical system with invariant measure $\nu$, admitting a spectral gap with decay exponent $\gamma$. Suppose $\Psi$ is a diffeomorphic delay embedding. Then for any smooth $f$-divergence,
\[
D_f\!\big(P^t\Gamma_K \,\Vert\, P^t\Psi(\Gamma)\big) \,\vee\, D_f\!\big(P^t\Psi(\Gamma)\,\Vert\, P^t\Gamma_K\big)
\;\le\; C e^{-\gamma t} + R_n,
\]
where $R_n$ is the finite-sample estimation error scaling with smoothness $s$, dimension $d_\mu$, and recurrence exponent $\rho$. In particular, as $n,t\to\infty$, both projections vanish, and the relative convergence rates diagnose whether the dynamics are tail-driven (slower e) or bulk-driven (slower m).
\end{theorem}

\paragraph{Remark.} These results establish the dictionary promised in the main text.

\section{Computational Routine for LIA}

This appendix provides a step-by-step description of how the least-information action (LIA) is evaluated in practice when the edge divergence is chosen as the Maximum Mean Discrepancy (MMD). The routine collapses the theoretical framework to an implementable sequence of computations.

\subsection{Primitive Setup}
\begin{enumerate}
    \item \textbf{Inputs.} A time series $\{y_t\}_{t=0}^n$, a kernel $k$ (RBF or Fisher), kernel scale $\ell$, and a candidate representation (embedding $h$, neighborhood metric/bandwidths).
    \item \textbf{Objective.} Estimate the \emph{edge action}
    \[
        A \;=\; \widehat{D}_{\mathrm{MMD}}^\nu \big( K_{\mathrm{data}}, K_{\mathrm{model}} \big)
        \;\approx\; \frac{1}{|\mathcal Q|} \sum_{Y_0 \in \mathcal Q} 
        \mathrm{MMD}^2\!\big( \widehat K_{\mathrm{data}}(\cdot|Y_0),\,\widehat K_{\mathrm{model}}(\cdot|Y_0) \big),
    \]
    where $\mathcal Q$ is a set of query points drawn from the test block, and report a block-bootstrap confidence interval (CI). If the CI contains $0$ and the mean is small, the empirical Rose condition is considered satisfied.
\end{enumerate}

\subsection{Local Mean Fit (for representation and geometry)}
\begin{enumerate}
    \item \textbf{Local linear mean.} For each query point $Y_0 \in \mathcal Q$, estimate
    \[
        Y_{t+1} \approx A(Y_0) Y_t + b(Y_0),
    \]
    using weighted least squares with exponential weights
    \[
        w_t = \exp\!\left(-\frac{\| \widetilde Y_t - \widetilde Y_0\|_2}{h}\right),
    \]
    where each coordinate of $Y$ is z-scored on the training block.
    \item \textbf{Residuals (for bootstrap).} Define $r_t = y_{t+1} - A(Y_0)^\top Y_t - b(Y_0)$; the weighted empirical residual law
    \(
      \widehat R_{Y_0} \propto \sum_t w_t \,\delta_{r_t}
    \)
    will be used to generate model-side fluctuations nonparametrically.
\end{enumerate}

\subsection{Divergence Estimation (MMD)}
For each $Y_0 \in \mathcal Q$, construct both empirical and model conditionals:

\begin{enumerate}
    \item \textbf{Empirical conditional.} Using the test block, form weights 
    \( \omega_t(Y_0) \propto \exp(-\|\widetilde Y_t-\widetilde Y_0\|_2/h_p) \)
    and define
    \(
      \widehat K_{\mathrm{data}}(\cdot|Y_0) = \sum_t \pi_t \,\delta_{y_{t+1}},\quad
      \pi_t = \omega_t / \sum_{t'} \omega_{t'}.
    \)

    \item \textbf{Model conditional (pushforward + residual bootstrap).} Draw $\epsilon \sim \widehat R_{Y_0}$ and set
    \(
      \tilde y_{t+1} = A(Y_0)^\top Y_0 + b(Y_0) + \epsilon,
    \)
    repeating $L$ times to obtain bootstrap samples from $\widehat K_{\mathrm{model}}(\cdot|Y_0)$.

    \item \textbf{MMD estimate.} With kernel $k_\ell(y,y')$, compute
    \[
    \mathrm{MMD}^2(\widehat P_{Y_0},\widehat Q_{Y_0})
    = \sum_{t,t'}\pi_t\pi_{t'}\,k_\ell(y_{t+1},y_{t'+1})
    + \frac{1}{L^2}\sum_{j,j'} k_\ell(\tilde y_j,\tilde y_{j'})
    - \frac{2}{L}\sum_{j}\sum_t \pi_t\,k_\ell(\tilde y_j,y_{t+1}).
    \]
    \item \textbf{Aggregate.} Average over $Y_0\in\mathcal Q$ to obtain $\widehat A$.
\end{enumerate}

\paragraph{Kernel choices.}
\begin{itemize}
  \item \textbf{RBF (default).} $k_\ell(y,y')=\exp(-\tfrac{(y-y')^2}{2\ell^2})$, with $\ell$ chosen by median heuristic or small validation sweep.
  \item \textbf{Fisher kernel (local FR geometry).} For local Gaussian mean model at $Y_0$,
  \[
    k_{Y_0}(y,y') = \frac{(y-\widehat\mu(Y_0))(y'-\widehat\mu(Y_0))}{\widehat\sigma(Y_0)^2},
  \]
  aligning MMD’s quadratic form with the Fisher–Rao metric.
\end{itemize}

\subsection{Information-Guided Basis and Lag Selection}
\begin{enumerate}
    \item \textbf{Pullback Fisher proxy.} Compute
    \[
        \bar{J} \;=\; \mathbb{E}_\nu[J(Y)] 
        \;\approx\; \frac{\sum_t w_t A_t^\top \big(\widehat{\mathrm{Var}}[r_t|Y_t]\big)^{-1} A_t}{\sum_t w_t}.
    \]
    \item \textbf{Basis selection.} For a linear projection $L \in \mathbb{R}^{d\times r}$ with $L^\top L = I_r$, maximize $\mathrm{trace}(L^\top \bar J L)$ (information PCA).
    \item \textbf{Lag selection.} Perform greedy D-optimal selection of lags by maximizing marginal gains in $\log \det (\bar J_{SS})$ via Schur complement.
\end{enumerate}

\subsection{Diagnostics}
\begin{itemize}
    \item \textbf{Learning curves.} Estimate $\hat A(n)$ versus $n$ on log--log scale, fit slope $\zeta$, and recover asymptotic error $A_\infty$.
    \item \textbf{Locality curves.} Estimate $\hat A(h)$ versus $h$ on log--log scale, fit slope $\eta$, and recover limiting error $A_0$.
    \item \textbf{Interpretation.} $\zeta,\eta$ identify effective smoothness $s$ and dimension $d$; $A_\infty \approx A_0 > 0$ signals an irreducible process-noise floor.
\end{itemize}

\section{Simulation Protocols}
\label{app:sim}

\subsection{Logistic Map with Tunable Diffusion}

\paragraph{Purpose.}
A minimal nonlinear, scalar system that demonstrates: near-zero edge action under good representation, recovery of learning/locality exponents $(\zeta,\eta)$, information-guided lag selection, and the emergence of an irreducible process-noise floor as diffusion increases.

\subsubsection{Data-Generating Model}
Let $x_t \in [0,1]$ evolve by
\[
x_{t+1} \;=\; r\,x_t\,(1-x_t) \;+\; \eta_t, 
\qquad r = 3.8,\quad \eta_t \sim \mathcal{N}(0,\sigma^2).
\]
Observations are $y_t := x_t$. Diffusion strength is swept over $\sigma \in \{0,\;0.005,\;0.01,\;0.02,\;0.05\}$.

\subsubsection{Protocol}
\begin{enumerate}
  \item \textbf{Trajectories.} For each $\sigma$, simulate $n=100{,}000$ points; discard a burn-in of $10{,}000$.
  \item \textbf{Splits.} Use contiguous $60\%/20\%/20\%$ train/val/test blocks.
  \item \textbf{Representation.} Delay-embed $y_t$ into $k$ lags: 
  $\;Y_t = [y_t, y_{t-1},\dots,y_{t-k+1}]^\top$, 
  with $k \in \{2,\dots,8\}$.
  \item \textbf{Local mean fit.}  
  For each $Y_0$ in the test block, fit $y_{t+1} \approx a(Y_0)^\top Y_t + b(Y_0)$ with exponential weights 
  \(
    w_t = \exp(-\|\widetilde Y_t-\widetilde Y_0\|_2/h).
  \)
  Record residuals $r_t$ and the weighted empirical residual law $\widehat R_{Y_0}$.
  \item \textbf{Hyperparameters.} Tune $h$ (and ridge) by validation conditional log-likelihood. Sweep $h$ on a log grid from $c_{\min}\cdot\mathrm{median}(d_t)$ to $c_{\max}\cdot\mathrm{median}(d_t)$ with $c_{\min}=0.25$, $c_{\max}=8$, using $m\in\{10,12,16\}$ points.
  \item \textbf{Edge divergence (MMD).}  
  For each test query $Y_0$:
  \begin{enumerate}
    \item Empirical: build $\widehat P_{Y_0}$ with test-block weights $h_p$.
    \item Model: bootstrap $L$ residual samples to form $\widehat Q_{Y_0}$.
    \item Compute $\mathrm{MMD}^2(\widehat P_{Y_0},\widehat Q_{Y_0})$ with kernel $k_\ell$.
  \end{enumerate}
  \item \textbf{Uncertainty.} Block bootstrap the test block with block length $B \approx 5\,\tau_{\mathrm{mix}}$ to produce a $95\%$ CI for $\widehat{A}$.
  \item \textbf{Information-guided lag selection.}  
  Estimate $\bar J$ on the training block. Perform greedy D-optimal lag selection over $\{1,\dots,8\}$; refit and recompute $\widehat{A}$.
  \item \textbf{Learning and locality exponents.}  
  (i) Compute $\widehat{A}(n)$ for subsample sizes $n\in\{10^3,2\!\times\!10^3,\dots,10^5\}$ and fit $\widehat{A}(n)\approx A_\infty + Cn_{\mathrm{eff}}^{-\zeta}$. \\
  (ii) Sweep $h$ across the log grid and fit $\widehat{A}(h)\approx A_0 + C h^\eta$.
\end{enumerate}

\subsubsection{Expected Recoveries}
\begin{itemize}
  \item \textbf{Edge action.} Adequate embeddings yield $\widehat{A}\approx 0$ with CI covering zero when $\sigma \approx 0$; misspecified embeddings yield $\widehat{A}>0$.
  \item \textbf{Exponents.} $\eta \approx 2s$ recovers smoothness $s$, while $\zeta \approx \tfrac{2s}{2s+d}$ identifies effective dimension $d$.
  \item \textbf{Noise floor.} As $\sigma$ grows, $A_\infty$ and $A_0$ coalesce above zero, revealing irreducible process noise.
  \item \textbf{Lag selection.} D-optimal lag choice stabilizes estimation, typically selecting $k\in\{3,4\}$ at $r{=}3.8$.
\end{itemize}

\subsubsection{Default Knobs}
\begin{itemize}
  \item Kernel: RBF with scale $\ell$ (median heuristic) or Fisher kernel.  
  \item Locality $h \in \{0.5,1,2,4\}\cdot n^{-1/5}$; decouple $(h_p,h_r)$ via multipliers $(\gamma_{hp},\gamma_{hr})$.  
  \item Model samples $L=5{,}000$.  
  \item Bootstrap resamples $=500$; block length $B \approx 5\,\tau_{\mathrm{mix}}$.  
\end{itemize}

\begin{algorithm}[H]
\caption{LIA with exponential kernel and conditional MMD}
\begin{algorithmic}[1]
\State \textbf{Input:} $\{y_t\}$, lag set $\mathcal K$, kernel $k$ (RBF or Fisher), resamples $B$
\State Split into Train/Val/Test; build delay embeddings for each $k\in\mathcal K$
\For{$h$ in log-grid,\; query $Y_0$ in Train/Val}
  \State Fit $y_{t+1}\approx a(Y_0)^\top Y_t+b(Y_0)$ by WLS with $w_t$; store residuals $r_t$
\EndFor
\State Choose $h$ by validation log-likelihood; set kernel scale $\ell$
\For{Test $Y_0$}
  \State \emph{Empirical} $\widehat P_{Y_0}$: weights $\omega_t \propto \exp(-\|\widetilde Y_t-\widetilde Y_0\|_2/h_p)$ on Test
  \State \emph{Model} $\widehat Q_{Y_0}$: sample $\epsilon \sim \widehat R_{Y_0}$; set $\tilde y = a(Y_0)^\top Y_0 + b(Y_0)+\epsilon$; repeat $L$ times
  \State Compute $\mathrm{MMD}^2(\widehat P_{Y_0},\widehat Q_{Y_0})$ with kernel $k_\ell$
\EndFor
\State $\widehat A \gets$ average MMD$^2$ over test queries; moving-block bootstrap ($B$ resamples) for CI
\State $\bar J \gets \tfrac{\sum w_t A_t^\top \widehat{\mathrm{Var}}[r_t|Y_t]^{-1} A_t}{\sum w_t}$; D-optimal lag selection $\to S^\star$; refit and recompute $\widehat A$ with CI
\State Build $\widehat A(n)$ and $\widehat A(h)$; fit slopes $(\zeta,\eta)$ and limits $(A_\infty,A_0)$
\State \textbf{Output:} $\widehat A$+CI, $S^\star$, $(\zeta,\eta)$, $A_\infty,A_0$, $\bar J$ spectrum and D-optimal gains
\end{algorithmic}
\end{algorithm}

\subsection{Numerical Demonstration Plan}

\subsubsection{Guiding Principle}
The theoretical results (Proposition~2 and Corollary~1 on bias--variance rates, Theorem~1 on measure conjugacy, and the new divergence--exponent bridge) require concrete simulation studies to validate the framework. The goal is to balance clarity with scope: provide simple, transparent evidence in this paper, while reserving extended stress tests and elaborations for a companion paper.

\subsubsection{Core Simulations (to include in this manuscript)}
\begin{enumerate}
    \item \textbf{Tier 1: One--dimensional controlled dynamics.} 
    \begin{itemize}
        \item System: autoregressive or scalar Markov process with conditional mean $\mu(y) = \mathrm{sign}(y)|y|^q + \eta_t$.
        \item Purpose: verify $\eta = 2s$ from locality curves and $\zeta = \tfrac{2s}{2s+d}$ from learning curves with $d=1$.
        \item Deliverables: 
            \begin{itemize}
                \item Log--log plot of $A(h)-A_0$ vs.\ $h$ (bias slope $\eta$).
                \item Learning curve $\min_h A_n(h)-A_0$ vs.\ $n$ (rate $\zeta$).
                \item Table: recovered $(\hat s, \hat d)$ compared to ground truth $(s,d)$.
            \end{itemize}
    \end{itemize}
    
    \item \textbf{Tier 3: Nonlinear 2D system with delay embedding.}
    \begin{itemize}
        \item System: H\'enon map with small noise; observation $Z_t = x_t$, embedded as $Y_t = [Z_t, Z_{t-1}, \dots, Z_{t-k+1}]$.
        \item Purpose: demonstrate Takens embedding + Rose equality $\Rightarrow$ measure conjugacy.
        \item Deliverables:
            \begin{itemize}
                \item Divergence (MMD--JS or JS) between empirical edge law $K_Y$ and pushforward $h_\#K_X$; report near-zero values (conjugacy test).
                \item Log--log locality and learning curves in embedded space, recovering effective $d$ smaller than ambient $k$.
                \item Witness heatmap for model misspecification (e.g.\ perturbed H\'enon parameter).
            \end{itemize}
    \end{itemize}

    \item \textbf{Tier 6: Dynamical exponents link.}
    \begin{itemize}
        \item \emph{Entropy/Lyapunov:} Logistic or H\'enon map, where true entropy rate and Lyapunov exponents are known. Purpose: check that divergence slopes align with known entropy growth and Lyapunov stretching.
        \item \emph{Mixing/Spectral:} Simple AR(1) or OU process, where correlation decay is exponential. Purpose: verify divergence contraction mirrors the mixing exponent (spectral gap).
        \item Deliverables:
            \begin{itemize}
                \item Divergence decay curves vs.\ $t$, compared to theoretical entropy/Lyapunov or mixing rates.
                \item Highlight slower e-projection in tail-driven regimes (recurrence), versus slower m-projection in bulk-driven regimes.
            \end{itemize}
    \end{itemize}
\end{enumerate}

\subsubsection{Extended Simulations (reserve for follow-up paper or appendix)}
\begin{enumerate}
    \item \textbf{Tier 2: State-dependent noise (heteroskedasticity).} Witness maps highlight local variance pockets invisible to RMSE.
    \item \textbf{Tier 4: Ornstein--Uhlenbeck with nonlinear drift.} Demonstrates path-action connection (Proposition~3), stability gains from quadratic action regularization.
    \item \textbf{Tier 5: Intrinsic vs.\ ambient dimension.} Synthetic manifold in high-dimensional space; estimator recovers effective dimension $d_\star$ rather than ambient dimension.
\end{enumerate}

\subsubsection{Summary}
\begin{itemize}
    \item \emph{This paper:} Tiers 1, 3, and 6 --- minimal but sufficient demonstrations linking theoretical rates, conjugacy, and dynamical exponents to concrete simulations.
    \item \emph{Companion work:} Tiers 2, 4, 5 --- richer stress tests and methodological elaborations, aimed at an applied ML or dynamical systems audience.
\end{itemize}


\end{document}
